\documentclass[11pt]{article}

\usepackage[margin=1in]{geometry}
\usepackage{amsmath}
\usepackage{amssymb}
\usepackage{booktabs}
\usepackage{siunitx}
\usepackage{graphicx}
\usepackage{hyperref}
\usepackage{longtable}
\usepackage{bm}

\title{Seawater Density Sensitivity to Temperature\
Using the TEOS-10 Gibbs SeaWater Oceanographic Toolbox}

\author{}
\date{\today}

\begin{document}

\maketitle

\begin{abstract}

This study examines the sensitivity of seawater density to a
$1,^\circ\mathrm{C}$ temperature perturbation using the
Thermodynamic Equation of Seawater 2010 (TEOS-10). The calculation
uses the Gibbs SeaWater (GSW) Oceanographic Toolbox directly from its
C source files, compiled and linked with a C++ application.

The input observations consist of in-situ temperature, Practical
Salinity, and depth. These observations are converted to sea pressure,
Absolute Salinity, and Conservative Temperature using TEOS-10
algorithms. Seawater density is subsequently calculated for the
original state and for a perturbed state in which temperature is
increased by $1,^\circ\mathrm{C}$.

The resulting density perturbation is calculated as the difference
between the perturbed and original densities. For the profiles examined
here, the temperature perturbation produces a decrease in seawater
density throughout the water column, with the magnitude of the response
varying as a function of depth and the thermodynamic state of the
seawater.

\end{abstract}

\section{Objective}

The objective of this calculation is to quantify the local
thermodynamic sensitivity of seawater density to temperature.

For each observation, the initial seawater state is represented by

\begin{equation}
\left(S_P,,t,,z\right),
\end{equation}

where

\begin{align}
S_P &= \text{Practical Salinity},\
t &= \text{in-situ temperature},\
z &= \text{vertical position}.
\end{align}

Depth is represented as a positive quantity, while TEOS-10 uses
height above the sea surface. Therefore,

\begin{equation}
z=-d,
\end{equation}

where $d$ is the measured depth below the sea surface.

The temperature perturbation is defined as

\begin{equation}
\Delta t = 1,^\circ\mathrm{C}.
\end{equation}

The perturbed temperature is consequently

\begin{equation}
t' = t+\Delta t.
\end{equation}

For this experiment,

\begin{equation}
\boxed{
t'=t+1,^\circ\mathrm{C}
}
\end{equation}

while salinity and pressure are held at their observed values.

The density perturbation is then defined as

\begin{equation}
\boxed{
\Delta\rho=\rho'-\rho
}
\end{equation}

where $\rho$ is the density of the original seawater state and
$\rho'$ is the density after the temperature perturbation.

\section{TEOS-10 Thermodynamic Framework}

The Thermodynamic Equation of Seawater 2010 (TEOS-10) describes the
thermodynamic properties of seawater using Absolute Salinity,
Conservative Temperature, and pressure.

The principal thermodynamic state variables used in this calculation
are therefore

\begin{equation}
\left(S_A,,CT,,p\right),
\end{equation}

where

\begin{align}
S_A &= \text{Absolute Salinity},\
CT &= \text{Conservative Temperature},\
p &= \text{sea pressure}.
\end{align}

Seawater density is represented as

\begin{equation}
\rho=\rho(S_A,CT,p).
\end{equation}

Equivalently, density may be expressed through specific volume $v$:

\begin{equation}
\boxed{
\rho=\frac{1}{v}
}
\end{equation}

with

\begin{equation}
v=v(S_A,CT,p).
\end{equation}

Thus,

\begin{equation}
\boxed{
\rho=
\frac{1}
{v(S_A,CT,p)}
}
\end{equation}

provides the fundamental thermodynamic relationship used by the
density calculation.

\section{Computational Sequence}

The complete calculation can be summarized as

\begin{equation}
\boxed{
(S_P,t,d,\phi,\lambda)
\longrightarrow
(p,S_A,CT)
\longrightarrow
\rho
}
\end{equation}

where $\phi$ and $\lambda$ denote latitude and longitude.

The perturbed calculation follows

\begin{equation}
\boxed{
(S_P,t+\Delta t,d,\phi,\lambda)
\longrightarrow
(p,S_A,CT')
\longrightarrow
\rho'
}
\end{equation}

with

\begin{equation}
\Delta t=1,^\circ\mathrm{C}.
\end{equation}

The final quantity of interest is

\begin{equation}
\boxed{
\Delta\rho
\rho(S_A,CT',p)
\rho(S_A,CT,p)
}
\end{equation}

where

\begin{equation}
CT'=CT(S_A,t+\Delta t,p).
\end{equation}

\section{Pressure From Depth}

The GSW function \texttt{gsw-p-from-z} is used to convert vertical
position into sea pressure.

The input depth is converted to TEOS-10 height according to

\begin{equation}
z=-d.
\end{equation}

The pressure is then evaluated as

\begin{equation}
p=
p(z,\phi,H_{\mathrm{dyn}}),
\end{equation}

where

\begin{align}
z &= \text{height relative to the sea surface},\
\phi &= \text{latitude},\
H_{\mathrm{dyn}} &= \text{dynamic height contribution}.
\end{align}

For the present calculation, the dynamic-height contribution is taken
according to the assumptions used by the implementation.

Pressure is reported in decibars:

\begin{equation}
p\quad [\mathrm{dbar}].
\end{equation}

Pressure is not assumed to be exactly proportional to depth. The
conversion depends upon gravitational acceleration and therefore
includes a latitude dependence.

\section{Practical Salinity and Absolute Salinity}

The observations provide Practical Salinity, $S_P$. TEOS-10 uses
Absolute Salinity, $S_A$, as its thermodynamic salinity variable.

The conversion is represented by

\begin{equation}
\boxed{
S_A =
S_A(S_P,p,\phi,\lambda)
}
\end{equation}

and is implemented using the GSW function

gsw-SA-from-SP
.

Absolute Salinity is expressed in

\begin{equation}
[\mathrm{g,kg^{-1}}].
\end{equation}

The distinction between $S_P$ and $S_A$ is important because the
TEOS-10 density calculation is based on Absolute Salinity rather than
Practical Salinity.

\section{Conservative Temperature}

The measured temperature is the in-situ temperature $t$.

TEOS-10 converts in-situ temperature to Conservative Temperature
according to

\begin{equation}
\boxed{
CT=CT(S_A,t,p)
}
\end{equation}

using the GSW function

gsw-CT-from-t
.

For the perturbed state,

\begin{equation}
t'=t+\Delta t,
\end{equation}

and therefore

\begin{equation}
\boxed{
CT'=CT(S_A,t+\Delta t,p)
}
\end{equation}

with

\begin{equation}
\Delta t=1,^\circ\mathrm{C}.
\end{equation}

An important feature of this calculation is that the perturbation is
applied to the original in-situ temperature before the TEOS-10
conversion to Conservative Temperature. Thus, the calculation does
not simply assume

\begin{equation}
CT'=CT+1,^\circ\mathrm{C}.
\end{equation}

Instead,

\begin{equation}
CT'
CT(S_A,t+1,^\circ\mathrm{C},p).
\end{equation}

\section{Density Calculation}

Once Absolute Salinity, Conservative Temperature, and pressure have
been determined, seawater density is calculated from

\begin{equation}
\boxed{
\rho=\rho(S_A,CT,p)
}
\end{equation}

and the perturbed density is

\begin{equation}
\boxed{
\rho'=\rho(S_A,CT',p).
}
\end{equation}

The density response is consequently

\begin{equation}
\boxed{
\Delta\rho
\rho(S_A,CT',p)
\rho(S_A,CT,p)
}
\end{equation}

with units

\begin{equation}
[\Delta\rho]
\mathrm{kg,m^{-3}}.
\end{equation}
\section{Local Temperature Sensitivity}

For a sufficiently small temperature perturbation, density may be
approximated using a first-order Taylor expansion:

\begin{equation}
\rho(T+\Delta T)
\approx
\rho(T)
+
\left(
\frac{\partial\rho}{\partial T}
\right)_{S_A,p}
\Delta T.
\end{equation}

Therefore,

\begin{equation}
\boxed{
\Delta\rho
\approx
\left(
\frac{\partial\rho}{\partial T}
\right)_{S_A,p}
\Delta T
}
\end{equation}

where the derivative is evaluated at constant Absolute Salinity and
pressure.

The thermal expansion coefficient $\alpha$ is defined by

\begin{equation}
\boxed{
\alpha
-\frac{1}{\rho}
\left(
\frac{\partial\rho}{\partial T}
\right)_{S_A,p}
}
\end{equation}

and therefore

\begin{equation}
\left(
\frac{\partial\rho}{\partial T}
\right)_{S_A,p}
-\rho\alpha.
\end{equation}

Substitution gives

\begin{equation}
\boxed{
\Delta\rho
\approx
-\rho\alpha\Delta T
}
\end{equation}

and, for the present experiment,

\begin{equation}
\boxed{
\Delta\rho
\approx
-\rho\alpha
}
\end{equation}

because

\begin{equation}
\Delta T=1,^\circ\mathrm{C}.
\end{equation}

The actual calculation does not rely on this linear approximation.
Instead, the complete TEOS-10 state is evaluated before and after the
finite temperature perturbation.

\section{Numerical Results}

The generated CSV file contains the following variables:

\begin{center}
\begin{tabular}{ll}
\toprule
Variable & Description \\
\midrule
\texttt{time} & Observation time \\
\texttt{depth\_m} & Depth below the sea surface \\
\texttt{temperature\_C} & Original in-situ temperature \\
\texttt{perturbed\_temperature\_C} & Temperature plus $1^\circ\mathrm{C}$ \\
\texttt{SP\_practical} & Practical Salinity \\
\texttt{SA\_absolute\_gkg} & Absolute Salinity \\
\texttt{pressure\_dbar} & Sea pressure \\
\texttt{CT\_C} & Original Conservative Temperature \\
\texttt{perturbed\_CT\_C} & Perturbed Conservative Temperature \\
\texttt{density\_kg\_m3} & Original density \\
\texttt{perturbed\_density\_kg\_m3} & Perturbed density \\
\texttt{density\_difference\_kg\_m3} & $\rho'-\rho$ \\
\bottomrule
\end{tabular}
\end{center}

Representative values from the generated profile are shown in
Table~\ref{tab:profile}.

\begin{table}[ht]
\centering
\caption{Representative TEOS-10 results from the generated profile.}
\label{tab:profile}
\begin{tabular}{rrrrr}
\toprule
Depth & $t$ & $S_A$ & $\rho$ & $\Delta\rho$ \\
(\si{m}) & ($^\circ\mathrm{C}$) & (\si{g\,kg^{-1}})
& (\si{kg,m^{-3}}) & (\si{kg,m^{-3}}) \\
\midrule
10 & 17.976 & 33.630 & 1024.1565 & -0.2484 \\
20 & 17.770 & 33.635 & 1024.2551 & -0.2468 \\
61 & 12.389 & 33.181 & 1025.2615 & -0.1968 \\
145 & 8.992 & 33.861 & 1026.7693 & -0.1657 \\
302 & 7.154 & 34.185 & 1028.0231 & -0.1504 \\
499 & 5.796 & 34.354 & 1029.2534 & -0.1406 \\
754 & 4.769 & 34.541 & 1030.7134 & -0.1358 \\
1001 & 3.974 & 34.634 & 1032.0298 & -0.1332 \\
2213 & 1.955 & 34.796 & 1037.9917 & -0.1416 \\
3873 & 1.509 & 34.863 & 1045.6019 & -0.1767 \\
\bottomrule
\end{tabular}
\end{table}


For example, at a depth of approximately \SI{10}{m}, the calculated
density is

\begin{equation}
\rho
1024.15649622\ \mathrm{kg,m^{-3}}.
\end{equation}

After increasing the temperature by $1,^\circ\mathrm{C}$,

\begin{equation}
\rho'
1023.90806198\ \mathrm{kg,m^{-3}}.
\end{equation}

The resulting density difference is

\begin{align}
\Delta\rho
&=
\rho'-\rho\
&=
1023.90806198
1024.15649622\
&=
-0.24843424\ \mathrm{kg,m^{-3}}.
\end{align}

Thus,

\begin{equation}
\boxed{
\Delta\rho
-0.24843424\ \mathrm{kg,m^{-3}}
}
\end{equation}

for this particular observation.

\section{Depth Dependence}

The magnitude of the density response changes with depth.

Near the surface, the profile produces density changes of approximately

\begin{equation}
\Delta\rho
\approx
-0.25\ \mathrm{kg,m^{-3}}
\end{equation}

for a $1,^\circ\mathrm{C}$ temperature increase.

Near \SI{1000}{m},

\begin{equation}
\Delta\rho
\approx
-0.13\ \mathrm{kg,m^{-3}}.
\end{equation}

At approximately \SI{3873}{m},

\begin{equation}
\Delta\rho
\approx
-0.177\ \mathrm{kg,m^{-3}}.
\end{equation}

This variation demonstrates that the temperature-density relationship
is not represented adequately by a single constant coefficient over the
entire water column.

The thermodynamic state of the seawater, including pressure,
temperature, and salinity, changes substantially with depth.

\section{Geographic Dependence}

Latitude is required by the TEOS-10 pressure calculation and may also
affect the conversion between Practical and Absolute Salinity.

The pressure calculation can therefore be represented schematically as

\begin{equation}
p=p(z,\phi,H_{\mathrm{dyn}}).
\end{equation}

The Absolute Salinity calculation depends upon geographic position:

\begin{equation}
S_A
S_A(S_P,p,\phi,\lambda).
\end{equation}

Consequently, latitude and longitude are part of the thermodynamic
calculation even though they are not themselves thermodynamic state
variables.

The current CSV output does not contain latitude or longitude fields.
Therefore, the CSV alone cannot be used to establish the geographic
location of the profile or its distance from the United States coast.

For complete reproducibility, the output should include

\begin{equation}
\phi=\text{latitude}
\end{equation}

and

\begin{equation}
\lambda=\text{longitude}.
\end{equation}

These quantities should preferably be recorded in the CSV alongside
each profile or in associated metadata.

\section{Software Implementation}

The calculation is implemented as a C++ program linked directly
against the required GSW C source files.

The principal source files are

\begin{itemize}
\item \texttt{cce1\_teos10\_density.cpp}
\item \texttt{gsw\_oceanographic\_toolbox.c}
\item \texttt{gsw\_saar.c}
\end{itemize}

along with the associated TEOS-10 header and data files.

The C source files are compiled into object files and linked directly
with the C++ application.

A representative compilation sequence is

\begin{verbatim}
gcc -c -O2 gsw_oceanographic_toolbox.c
gcc -c -O2 gsw_saar.c

g++ -O2 cce1_teos10_density.cpp
gsw_oceanographic_toolbox.o
gsw_saar.o
-o cce1_teos10_density
\end{verbatim}

The resulting executable produces

\begin{verbatim}
cce1_teos10_density.csv
\end{verbatim}

containing the original and perturbed thermodynamic states.

\section{Interpretation}

The results demonstrate that increasing seawater temperature by
$1,^\circ\mathrm{C}$ decreases seawater density for the states
represented by the profile.

The response is not constant with depth because seawater density is a
function of multiple thermodynamic variables:

\begin{equation}
\rho=\rho(S_A,CT,p).
\end{equation}

The experiment isolates the temperature component by holding the
observed salinity and pressure fixed while modifying temperature.

Thus, the experiment should be interpreted as a controlled
thermodynamic sensitivity calculation rather than as a complete model
of future ocean warming.

In particular, it does not model simultaneous changes in salinity,
mixing, circulation, evaporation, precipitation, air-sea heat flux,
or water-mass formation.

\section{Limitations}

The principal limitation is that the experiment applies a prescribed
temperature perturbation while holding other measured quantities
fixed.

The calculation therefore evaluates

\begin{equation}
\rho(S_A,CT+\Delta CT,p)
\rho(S_A,CT,p)
\end{equation}

under the prescribed perturbation rather than simulating the full
physical evolution of the ocean.

A further limitation is that the current CSV file does not contain
latitude and longitude. These values should be retained as metadata
for future analyses because they affect the pressure and Absolute
Salinity calculations.

\section{References}

\begin{enumerate}

\item International Oceanographic Commission, Scientific Committee on
Oceanic Research, and International Association for the Physical
Sciences of the Oceans.
\textit{The International Thermodynamic Equation of Seawater---2010:
Calculation and Use of Thermodynamic Properties}.
TEOS-10 Manual, UNESCO, 2010.

\item McDougall, T. J., Jackett, D. R., Wright, D. G., and Feistel, R.
``Accurate and Computationally Efficient Algorithms for Potential
Temperature and Density of Seawater.''
\textit{Journal of Atmospheric and Oceanic Technology},
Vol. 20, 2003, pp. 730--741.

\item Roquet, F., Madec, G., McDougall, T. J., and Barker, P. M.
``Accurate Polynomial Expressions for the Density and Specific Volume
of Seawater Using the TEOS-10 Standard.''
\textit{Ocean Modelling}, Vol. 90, 2015, pp. 29--43.

\item Saunders, P. M.
``Practical Conversion of Pressure to Depth.''
\textit{Journal of Physical Oceanography}, Vol. 11, 1981.

\item TEOS-10.
\textit{Gibbs SeaWater (GSW) Oceanographic Toolbox}.
Available at:
\url{https://www.teos-10.org/pubs/gsw/html/gsw_contents.html}

\item TEOS-10.
\textit{TEOS-10 Manual}.
Available at:
\url{https://www.teos-10.org/pubs/TEOS-10_Manual.pdf}

\end{enumerate}

\end{document}
